\noindent\textbf{Por qué se usa este procedimiento.} La idealización concentra área axial en cuatro larguerillos: $I_z=\sum B_i y_i^2=4By^2$. Al cruzar un larguerillo, el primer momento cambia en $By$, de modo que el salto de flujo es $\Delta q=C_yBy/I_z$.

\begin{enumerate}

\item La inercia es $4 B y^2$.

\end{enumerate}

\noindent\textbf{Operaciones y resultados intermedios.}

\begin{enumerate}

\item Salto de flujo.

\[\fitmath{y_{1}=\frac{6000\cdot 300\cdot 80}{4\cdot 300\cdot {80}^{2}}\approx 18.75\,\mathrm{N/mm}}\]

\end{enumerate}
